% =====================================================================
%  3 Trabalhos Relacionados
% =====================================================================
\section{Trabalhos Relacionados}
\label{sec:relacionados}

Os trabalhos relacionados a este estudo se organizam em quatro linhas: a
criação de rotinas em assistentes de voz, as dinâmicas entre moradores em
casas inteligentes, a avaliação de comunicabilidade com a MIS e a análise
de avaliações de usuários. A Tabela~\ref{tab:relacionados} compara os
principais trabalhos com a proposta deste artigo.

\begin{table}[!htbp]
  \centering
  \caption{Comparação entre os trabalhos relacionados e este artigo.}
  \label{tab:relacionados}
  \small
  \begin{tabular}{>{\raggedright\arraybackslash}p{3.2cm}
                  >{\raggedright\arraybackslash}p{5.6cm}
                  >{\raggedright\arraybackslash}p{2.1cm}
                  >{\raggedright\arraybackslash}p{2.1cm}}
    \hline
    \textbf{Trabalho} & \textbf{Objeto e método} &
    \textbf{Múltiplos moradores} & \textbf{Participantes} \\
    \hline
    \citet{barricelli2024routines} & Criação de rotinas por voz e multimodal; experimentos controlados & Não & Sim \\
    \citet{king2024sasha} & Rotinas geradas por LLM; levantamento e estudo com usuários & Não & Sim \\
    \citet{geeng2019control} & Relações entre moradores; entrevistas e amostragem de experiência & Sim & Sim \\
    \citet{koshy2021passenger} & Pilotos e passageiros, incluindo rotinas; entrevistas e levantamento & Sim & Sim \\
    \citet{zeng2019multiuser} & Controle de acesso multiusuário; protótipo implantado em domicílios & Sim & Sim \\
    \citet{carvalho2023simsr} & Comunicabilidade e acessibilidade; SIM-SR & Não & Não \\
    \citet{atefi2020reviews} & Problemas em \textit{skills} da Alexa; análise de avaliações & Não & Não \\
    \hline
    Este trabalho & Rotinas no aplicativo Alexa; MIS com lente para múltiplos interlocutores e avaliações públicas & Sim & Não \\
    \hline
  \end{tabular}
\end{table}

\subsection{Criação de rotinas em assistentes de voz}

\citet{barricelli2022multimodal} propuseram criar rotinas usando os
recursos multimodais (visão, fala e toque) da Alexa em dispositivos Echo
Show e relataram que os usuários consideraram essa forma mais fácil do que
a interação habitual com o aplicativo Alexa. Em seguida,
\citet{barricelli2024routines} desenvolveram \textit{skills} da Amazon e
\textit{actions} do Google para transferir a criação de rotinas do
aplicativo para o próprio alto-falante. Em dois experimentos controlados
com 40 participantes, a interação multimodal foi preferida à interação
apenas por voz, e a qualidade percebida pareceu depender da marca do
dispositivo.

Com os grandes modelos de linguagem (LLMs), surgiram propostas de criar
rotinas a partir de pedidos em linguagem natural. O Sasha responde a
comandos pouco especificados, como ``deixe aconchegante'', executando
planos com os dispositivos disponíveis ou criando rotinas de automação;
seus autores mostram que os LLMs raciocinam de forma criativa, mas
apresentam padrões de falha que reduzem sua
utilidade~\cite{king2024sasha}. O GreenIFTTT usa um LLM para ajudar o
usuário a criar rotinas de economia de energia e foi avaliado quanto à
usabilidade em um estudo exploratório com 13
participantes~\cite{giudici2024greenifttt}.

Esses trabalhos buscam facilitar a autoria de rotinas e avaliam a
experiência de quem as cria. Nenhum deles analisa o que a interface de
criação comunica sobre as demais pessoas da casa, que são afetadas pelas
rotinas sem participar de sua criação.

\subsection{Dinâmicas entre moradores em casas inteligentes}

Os estudos apresentados na Seção~\ref{sec:multiusuario} investigam as
relações entre moradores por meio de entrevistas, amostragem de
experiência, levantamentos e implantações em domicílios
reais~\cite{geeng2019control,koshy2021passenger,zeng2017concerns,zeng2019multiuser}.
Deles derivam recomendações de design que dizem respeito diretamente às
interfaces de configuração. \citet{geeng2019control} recomendam que, ao
configurar um dispositivo, os aplicativos perguntem se há outras pessoas
na casa e facilitem o acesso delas. \citet{koshy2021passenger} propõem
controles que aproximem os passageiros das funcionalidades adequadas aos
seus hábitos e que lhes deem mais controle sobre os dados que
compartilham. \citet{zeng2019multiuser} prototiparam um aplicativo com
controle de acesso e notificações de atividade e observaram que, em
algumas casas, as normas sociais de convivência dispensaram os controles
técnicos, enquanto o aumento da complexidade das opções reduziu a
usabilidade.

Esses trabalhos partem da experiência relatada pelos moradores. A forma
como as interfaces comerciais comunicam a autoria, o controle e os
efeitos das rotinas, que pode ser examinada sem envolver participantes,
não é objeto de análise sistemática nesses estudos.

\subsection{Avaliação de comunicabilidade com a MIS}

A MIS tem sido aplicada e adaptada a diferentes contextos. A MIS
Científica foi ilustrada por seus autores na análise de como o editor
Simple CSS comunica seu \textit{feedback} a usuários não especialistas,
com resultados triangulados com material publicado na web sobre o
sistema, incluindo opiniões de usuários, e com a inspeção de outro
sistema~\cite{desouza2010scientific,fortunato2025ihc}. A SIM-SR foi
aplicada em três estudos de caso e levou os avaliadores a considerar o
impacto dos problemas de acessibilidade na comunicabilidade, indo além da
identificação de violações das diretrizes WCAG~\cite{carvalho2023simsr}.
\citet{fortunato2025ihc} sistematizam essas e outras variantes e indicam
que a combinação de versões do método em um mesmo sistema ainda precisa
ser investigada.

Até onde pudemos verificar, a MIS não foi aplicada a aplicativos de
configuração de assistentes de voz nem a rotinas de casas inteligentes.

\subsection{Análise de avaliações de usuários}

As avaliações publicadas por usuários em lojas de aplicativos são uma
fonte de evidência que não exige recrutamento. \citet{atefi2020reviews}
analisaram mais de 100 mil avaliações de mais de 2.800 \textit{skills} da
Alexa e, por meio de codificação aberta, identificaram 16 tipos de
problemas, dos quais os mais frequentes envolvem conteúdo, integração com
serviços e dispositivos, erros e regressões. O estudo trata das
\textit{skills}, e não do aplicativo usado para configurar o assistente e
suas rotinas.

\subsection{Posicionamento deste trabalho}

A Tabela~\ref{tab:relacionados} resume os trabalhos discutidos. Este
artigo combina três aspectos que aparecem separados na literatura: o foco
nas rotinas, a consideração de múltiplos moradores e uma avaliação que não
envolve participantes. Para isso, aplica a MIS com uma lente para
múltiplos interlocutores e triangula os resultados com avaliações
públicas do aplicativo, seguindo a estrutura da MIS
Científica~\cite{desouza2010scientific}.
